\documentclass[12pt]{article}
\usepackage[margin=1in]{geometry}
\usepackage{amsmath,amssymb}
\usepackage{graphicx}
\usepackage{booktabs}
\usepackage{hyperref}
\usepackage{natbib}
\usepackage{setspace}
\usepackage{lineno}
\doublespacing
\linenumbers

\title{Lateral Entorhinal Cortex Neurons Encode Distinct Experiential Epochs Around Reward During Goal-Directed Navigation in Mice}

\author{Author One$^{1,*}$, Author Two$^{1}$, Author Three$^{2}$ \\
\small $^1$Princeton Neuroscience Institute, Princeton University, Princeton, NJ, USA \\
\small $^2$Department of Molecular Biology, Princeton University, Princeton, NJ, USA \\
\small $^*$Corresponding author: author.one@princeton.edu}

\date{}

\begin{document}
\maketitle

\begin{abstract}
The hippocampus integrates spatial (``where'') and experiential (``what'') information to support episodic memory and flexible navigation, but the source of the experiential signal is poorly understood. Here we developed a microprism-based two-photon imaging approach to record from the lateral entorhinal cortex (LEC) in head-fixed mice navigating a virtual linear track for water reward. We identified three functionally distinct neuronal populations in LEC: pre-reward neurons encoding goal approach, reward-consumption-active (RCA) neurons, and post-reward neurons encoding departure. When the reward location was moved, pre-reward and post-reward populations immediately shifted their activity to track the new reward site while maintaining their functional identity. These populations were largely invariant to spatial location and environment, distinguishing them from medial entorhinal cortex (MEC) spatial representations. A hidden Markov model analysis revealed that pre-reward population activity is better described by abrupt state transitions at variable locations than by gradual recruitment. Optogenetic inhibition of LEC disrupted learning of new reward locations but did not impair behavior at familiar locations, establishing a causal role for LEC in reward-location learning. These findings demonstrate that LEC provides the non-spatial experiential context signals that the hippocampus requires for building reward-centric cognitive maps.
\end{abstract}

\section{Introduction}

The hippocampus has long been recognized as a critical structure for spatial navigation and episodic memory \citep{okeefe1971hippocampus}. Place cells in the hippocampal area CA1 encode the animal's current location in an environment, and their activity is modulated by contextual factors including the presence and location of reward \citep{dupret2010reorganization, gauthier2018dedicated}. A central question in systems neuroscience is how spatial information is integrated with non-spatial experiential information---what happened at a particular location, and when---to form the rich episodic memories that support flexible behavior.

The entorhinal cortex provides the principal cortical input to the hippocampus and has been divided into medial and lateral subdivisions with distinct functional roles. The medial entorhinal cortex (MEC) is home to grid cells, border cells, and head-direction cells that provide the spatial metric for hippocampal maps \citep{hafting2005microstructure, stensola2012entorhinal}. In contrast, the lateral entorhinal cortex (LEC) has been proposed to carry non-spatial information about objects, contexts, and temporal sequences \citep{deshmukh2011representation, tsao2018traces}. However, the specific contribution of LEC to reward-guided navigation has remained poorly understood, in part because the anatomical position of LEC makes it challenging to access with standard imaging approaches.

During goal-directed navigation, hippocampal place cells show characteristic reward-related modulation: overrepresentation of place fields near reward locations and distinct firing patterns during goal approach versus departure \citep{gauthier2018dedicated, eichenbaum2014time}. These observations imply that the hippocampus receives an upstream signal indicating the animal's experiential state relative to reward---approaching, consuming, or departing. The MEC, with its predominantly spatial coding, is unlikely to provide this signal. The LEC, by contrast, is well positioned to encode non-spatial experiential context \citep{woods2020history}, but direct evidence for reward-experience-epoch coding in LEC during navigation has been lacking.

In this study, we overcame the technical challenge of imaging LEC in behaving mice by developing a novel microprism-based two-photon imaging preparation. We recorded calcium transients from approximately 500 active neurons per imaging field across 47 imaging fields in GCaMP6s transgenic mice \citep{chen2013ultrasensitive}. We discovered three functionally distinct populations in LEC that encode the approach, consumption, and departure phases of goal-directed navigation, and we characterized their response properties across reward relocations, environment changes, and optogenetic manipulations.

\section{Methods}

\subsection{Animals and Surgical Preparation}

All experiments were approved by the institutional animal care and use committee. Transgenic mice expressing GCaMP6s were used for two-photon imaging experiments. A separate cohort of mice expressing Arch in LEC neurons was used for optogenetic experiments. Mice were implanted with a custom cranial window consisting of a 3 mm No.~0 glass coverslip with an attached 2.0 $\times$ 2.0 mm square microprism (45-degree angle). The microprism rotated the imaging plane by 90 degrees, allowing optical access to the laterally oriented LEC through a conventional upright two-photon microscope at depths exceeding 250 $\mu$m. Fields of view were approximately 700 $\times$ 700 $\mu$m.

\subsection{Virtual Reality Behavioral Paradigm}

Head-fixed mice ran on a treadmill through a 3.1 m linear virtual reality track. Water reward was delivered at a specific location on the track (standard location: 2.3 m; alternative locations: 0.7 m, 1.5 m). The track was visually cue-rich, but the reward location was not marked by any visual feature. Session duration was approximately 40 minutes. Mice were considered to have learned the task when they completed at least 2 laps per minute with anticipatory licking behavior. Velocity was binned at 1 cm intervals, with stationary periods excluded from analysis.

\subsection{Two-Photon Imaging and Cell Extraction}

Imaging was performed with a conventional upright two-photon microscope. Motion correction was applied using Suite2p \citep{pachitariu2017suite2p}. A novel iterative algorithm was used for neuropil correction, baseline adjustment, and deconvolution of calcium transients. Across all sessions, 47 imaging fields yielded a mean of 496 active neurons per field (range: 150--843).

\subsection{Spatial Cell Classification}

Spatial cells were identified by comparing the distribution of firing peaks across track positions to shuffled data. Cells were classified as pre-reward if their firing peak occurred before the reward location, post-reward if their peak occurred after the reward location, and reward-consumption-active (RCA) if they showed peak activity within the first second following reward delivery.

\subsection{Reward Relocation Protocol}

In reward relocation experiments, the reward was moved from the familiar location to a new location mid-session. Approximately 15 minutes of data were collected at the familiar location, followed by approximately 15 minutes at the new location.

\subsection{Hidden Markov Model Analysis}

To distinguish between recruitment and state-transition models of pre-reward population activation, a two-state hidden Markov model (HMM) was fit to population activity in the 10-second window before reward delivery. The two states were ``inactive'' and ``active'' (absorbing), and the HMM learned emission and transition probabilities from the data.

\subsection{Optogenetic Inhibition}

Arch-expressing mice were implanted with bilateral optical fibers over LEC. In alternating blocks, light was delivered (Light ON) or withheld (Light OFF) during task performance. The primary behavioral readout was the learning rate for a newly introduced reward location compared to performance at the familiar location.

\section{Results}

\subsection{Prominent Reward-Related Firing in LEC}

Well-trained mice exhibited anticipatory deceleration and licking behavior before reaching the reward location. Two-photon imaging revealed that LEC contained a large proportion of neurons with spatially modulated firing, with a prominent enrichment of firing fields near the reward location. This enrichment was characterized by a strong bimodal clustering of spatial cell peaks before and after reward, distinct from the more uniform distribution observed in MEC and the qualitatively different reward-related overrepresentation seen in CA1 (Table~\ref{tab:regions}).

\begin{table}[ht]
\centering
\caption{Comparison of reward-related neural properties across brain regions.}
\label{tab:regions}
\begin{tabular}{lcccc}
\toprule
Feature & LEC & MEC & CA1 \\
\midrule
Spatial cells & Yes & Yes & Yes \\
Reward clustering & Prominent & Minimal & Present \\
Pre-reward population & Distinct & Not prominent & Present (different) \\
Post-reward population & Distinct & Not prominent & Present (different) \\
RCA neurons & Yes & Minimal & Yes \\
Response to reward move & Immediate shift & Little change & Some redistribution \\
Environment invariance & Yes & No & No \\
\bottomrule
\end{tabular}
\end{table}

\subsection{Three Functionally Distinct Populations}

Analysis of LEC firing patterns revealed three clearly separable populations: (1) pre-reward neurons that were active during the approach to the reward location, (2) reward-consumption-active (RCA) neurons that fired during reward delivery and consumption, and (3) post-reward neurons that were active during departure from the reward site. These populations showed minimal overlap, indicating functional specialization within LEC for encoding distinct experiential epochs around reward.

\subsection{Immediate Tracking of Reward Relocation}

When the reward location was moved mid-session (e.g., from 2.3 m to 1.5 m), all three populations immediately shifted their activity to track the new reward location (Table~\ref{tab:relocation}). Pre-reward neurons that had been active before the old reward location began firing before the new location within the first few laps. Post-reward neurons similarly shifted to fire after the new location. RCA neurons shifted to fire during consumption at the new site. Critically, neurons maintained their functional identity---pre-reward neurons remained pre-reward neurons, post-reward remained post-reward---demonstrating that these populations encode the experiential epoch relative to reward rather than absolute spatial position.

\begin{table}[ht]
\centering
\caption{Population responses to reward relocation.}
\label{tab:relocation}
\begin{tabular}{lll}
\toprule
Population & Response to Reward Move & Latency \\
\midrule
Pre-reward & Shifts to new approach zone & Within first few laps \\
Post-reward & Shifts to new departure zone & Within first few laps \\
RCA & Shifts to new consumption site & Immediate \\
\bottomrule
\end{tabular}
\end{table}

In contrast, MEC neurons showed minimal shift in their firing patterns following reward relocation, consistent with their role in encoding allocentric spatial information rather than reward-relative experience. CA1 neurons showed some redistribution of place fields, reflecting the integration of both spatial and experiential inputs.

\subsection{Environment Invariance}

When mice were switched between different virtual environments (environment A to environment B), the pre-reward and post-reward populations maintained their functional identity. The same neurons that were active before reward in environment A were active before reward in environment B, and likewise for post-reward neurons. This environment invariance contrasts sharply with MEC and CA1, where neural representations remap across environments, and demonstrates that LEC reward-experience coding is abstracted from specific environmental features.

\subsection{Pre-Reward Activity Reflects State Transitions}

At the population mean level, pre-reward neuron activity appeared as a gradual ramp-up prior to reward delivery. However, HMM analysis of single-trial activity revealed that this apparent ramp was better explained by abrupt state transitions occurring at variable locations across laps (Table~\ref{tab:hmm}). The recruitment model, which posits gradually increasing numbers of active cells, provided a poor fit to single-trial data. In contrast, the state-transition model, in which the entire pre-reward population activates together at a variable location, better captured the binary nature of activation observed on individual laps.

\begin{table}[ht]
\centering
\caption{Hidden Markov model analysis of pre-reward population dynamics.}
\label{tab:hmm}
\begin{tabular}{lll}
\toprule
Model & Description & Fit Quality \\
\midrule
Recruitment & Gradual increase in active cells per trial & Poor (single trial) \\
State transition & Population activates at variable location & Better fit \\
\midrule
\multicolumn{3}{l}{HMM parameters: 2 states (inactive, active/absorbing)} \\
\multicolumn{3}{l}{Time window: 10 s before reward delivery} \\
\bottomrule
\end{tabular}
\end{table}

\subsection{RCA Neuron Properties}

RCA neurons in LEC showed several distinctive properties. Their activity was temporally locked to reward delivery time, with peak firing occurring within the first second after reward. They were spatially invariant, firing at the reward location regardless of where on the track reward was positioned. They were environment-invariant, maintaining their response across different virtual environments. And they immediately tracked reward relocation, confirming their role in encoding the consumption experience rather than any particular spatial location.

\subsection{Optogenetic Inhibition Disrupts Learning}

Optogenetic inhibition of LEC via bilateral Arch activation significantly disrupted mice's ability to learn a new reward location. Under Light ON conditions, mice showed slowed acquisition of anticipatory licking and deceleration at the new reward site compared to Light OFF controls. However, LEC inhibition did not impair behavior at an already-learned reward location, indicating that LEC is specifically required for new reward-location learning rather than for the expression of previously acquired reward associations.

\subsection{Layer-Specific Analysis}

Analysis of reward clustering as a function of imaging depth within LEC revealed that both superficial and deep layers contained reward-related firing, though with some differences in magnitude. ANOVA testing the effects of imaging depth (50 $\mu$m bins), anticipatory deceleration, sex, and age on reward clustering confirmed that the reward-enrichment phenomenon was robust across these factors.

\section{Discussion}

Our findings demonstrate that the lateral entorhinal cortex contains three functionally distinct neuronal populations that encode the experiential epochs of goal-directed navigation: approach to reward, consumption, and departure. These populations are reward-centric rather than spatially anchored, shifting immediately when reward location changes while maintaining their functional identity. This organization positions LEC as a source of the non-spatial ``what'' information that the hippocampus requires for contextualizing spatial maps during goal-directed behavior \citep{rabinowitz2012constructing}.

The immediate tracking of reward relocation by all three populations, coupled with their environment invariance, distinguishes LEC representations from those in MEC and CA1. MEC neurons encode space in an allocentric, environment-specific manner \citep{hafting2005microstructure}, while CA1 place cells show complex remapping that reflects the integration of spatial and contextual inputs \citep{okeefe1971hippocampus}. LEC provides a parallel, stable representation of the experiential context that is invariant to the specific environment but flexibly updates to track changes in reward contingencies.

The HMM analysis of pre-reward population dynamics resolves an apparent discrepancy between single-trial and trial-averaged views of pre-reward activity. The gradual ramp observed at the mean level arises from averaging across trials in which the state transition occurs at different locations. This variability in transition timing may reflect trial-to-trial fluctuations in the animal's internal state, attentional engagement, or running strategy.

The optogenetic finding that LEC inhibition disrupts new reward-location learning but spares behavior at familiar locations indicates a specific role for LEC in encoding novel reward experiences. Once a reward association has been established, the hippocampal system can maintain the representation independently of ongoing LEC input. This distinction has implications for understanding the dynamic interaction between cortical and hippocampal memory systems during learning.

The microprism imaging approach developed here provides a general solution for accessing laterally oriented brain structures in behaving mice. The ability to record from approximately 500 neurons per field of view across extended behavioral sessions opens new possibilities for studying LEC function in a variety of behavioral paradigms beyond the reward-based navigation studied here.

\subsection{Limitations}

Several limitations merit consideration. First, all recordings were performed in head-fixed mice navigating a one-dimensional virtual environment, which constrains the spatial complexity of the task compared to free navigation. Second, the calcium imaging approach, while providing excellent spatial coverage, has lower temporal resolution than electrophysiology, potentially blurring rapid dynamics. Third, the optogenetic inhibition affected all LEC neurons rather than specific subpopulations, leaving open the question of which population is most critical for learning.

\section{Conclusion}

We identify three dedicated neuronal populations in the lateral entorhinal cortex that encode the experiential epochs of goal-directed navigation relative to reward. These populations are reward-centric, environment-invariant, and causally required for learning new reward locations. The parallel encoding of experiential context in LEC alongside spatial context in MEC provides a dual-input framework for understanding how the hippocampus constructs the rich, event-structured maps that support episodic memory and flexible navigation.

\section*{Acknowledgments}
We thank all members of the laboratory for helpful discussions and technical assistance.

\bibliographystyle{plainnat}
\bibliography{references}

\end{document}
